% Chapter 2, Topic _Linear Algebra_ Jim Hefferon
%  http://joshua.smcvt.edu/linearalgebra
%  2001-Jun-11
\topic{Analisis Dimensi}
``Anda tidak dapat menjumlahkan apel dan jeruk,'' demikian bunyi pepatah lama.
Pepatah itu mencerminkan pengalaman kita bahwa dalam penerapan, besaran-besaran
memiliki satuan dan pencatatan satuan tersebut dapat membantu.
Semua orang pernah melakukan perhitungan seperti berikut ini,
yang menggunakan satuan untuk memeriksa hasilnya.
\begin{equation*}
  60\,\frac{\text{detik}}{\text{menit}}
  \cdot 60\,\frac{\text{menit}}{\text{jam}}
  \cdot 24\,\frac{\text{jam}}{\text{hari}}
  \cdot 365\,\frac{\text{hari}}{\text{tahun}}
  = 31\,536\,000\,\frac{\text{detik}}{\text{tahun}}
\end{equation*}
Kita dapat mengembangkan gagasan mencantumkan satuan melampaui sekadar pencatatan.
Kita dapat menggunakan satuan untuk menarik kesimpulan
tentang hubungan apa saja yang mungkin di antara besaran-besaran fisik.

Sebagai permulaan, perhatikan persamaan benda jatuh
$\text{jarak}=16\cdot(\text{waktu})^2$.
Jika jarak dinyatakan dalam kaki dan waktu dalam detik, maka pernyataan ini
benar.
Namun, persamaan itu tidak benar dalam sistem satuan lain, seperti meter dan detik,
karena~$16$ bukan konstanta yang
tepat dalam sistem tersebut.
Kita dapat memperbaikinya dengan menambahkan satuan pada $16$, sehingga menjadi
\definend{konstanta berdimensi}\index{dimensi!konstanta}.
\begin{equation*}
  \text{jarak}=16\,\frac{\text{kaki}}{\text{detik}^2}\cdot (\text{waktu})^2
\end{equation*}  
Sekarang persamaan itu juga berlaku dalam sistem meter-detik
karena ketika kita menyelaraskan satuannya (satu kaki kira-kira $0.30$ meter),
\begin{equation*}
  \text{jarak dalam meter}=16\,\frac{0.30\,\text{m}}{\text{detik}^2}
                    \cdot (\text{waktu dalam detik})^2
                   =4.8\,\frac{\text{m}}{\text{detik}^2}
                    \cdot (\text{waktu dalam detik})^2
\end{equation*}
konstantanya ikut disesuaikan.
Jadi, agar memperoleh persamaan yang benar dalam berbagai sistem satuan,
kita membatasi perhatian
pada persamaan yang menggunakan konstanta berdimensi.
Persamaan semacam itu disebut \definend{lengkap}.\index{persamaan lengkap}

Dengan melepaskan diri dari suatu sistem satuan tertentu, kita cukup
mengukur besaran sebagai kombinasi
satuan panjang~$L$, massa~$M$, dan waktu~$T$.
Ketiganya merupakan
\definend{dimensi fisik}\index{dimensi!fisik} kita.
Sebagai contoh, kita dapat mengukur kecepatan
dalam $\text{kaki}/\text{detik}$
atau $\text{depa}/\text{jam}$, tetapi bagaimanapun juga besaran itu melibatkan
satuan panjang dibagi satuan waktu,
sehingga \definend{rumus dimensi}\index{dimensi!rumus}
kecepatan adalah $L/T$.
Demikian pula, rumus dimensi massa jenis adalah $M/L^3$.

Untuk menuliskan rumus dimensi,
kita akan menggunakan eksponen negatif sebagai pengganti
pecahan, dan akan mencantumkan dimensi yang memiliki
eksponen nol.
Dengan demikian, kita akan menuliskan rumus dimensi
kecepatan sebagai $L^1M^0T^{-1}$ dan rumus dimensi massa jenis
sebagai $L^{-3}M^1T^0$.

Dengan itu, ``Anda tidak dapat menjumlahkan apel dan jeruk'' menjadi
anjuran untuk memeriksa bahwa semua suku dalam suatu persamaan memiliki
rumus dimensi yang sama.
Salah satu contohnya adalah bentuk persamaan benda jatuh berikut,~$d-gt^2=0$.
Rumus dimensi suku $d$ adalah $L^1M^0T^0$.
Untuk suku lainnya, rumus dimensi $g$ adalah $L^1M^0T^{-2}$
($g$ di atas diberikan sebagai $16\,\text{kaki}/\text{detik}^2$),
dan rumus dimensi $t$ adalah $L^0M^0T^1$, sehingga rumus dimensi
seluruh suku $gt^2$ adalah
$L^1M^0T^{-2}(L^0M^0T^1)^2=L^1M^0T^0$. 
Jadi, kedua suku tersebut memiliki rumus dimensi yang sama.
Persamaan dengan sifat ini disebut \definend{homogen secara dimensi}.

Besaran dengan rumus dimensi $L^0M^0T^0$ disebut \definend{tak berdimensi}.
Sebagai contoh, kita mengukur sebuah sudut dengan mengambil
rasio panjang busur yang dibentuk oleh sudut tersebut terhadap jari-jari,
\begin{center}
  \includegraphics{vs/mp/ch2.22}
\end{center}
yang merupakan rasio suatu panjang terhadap suatu panjang,
$(L^1M^0T^0)(L^1M^0T^0)^{-1}$, sehingga sudut memiliki rumus dimensi
$L^0M^0T^0$.
%(Of course, to measure the angle in degrees instead we would 
%multiply by a scalar, which would leave the dimensional formula unchanged.)

Contoh klasik penggunaan satuan untuk lebih dari sekadar pencatatan, yakni
untuk menarik kesimpulan, meninjau rumus periode sebuah pendulum.
\begin{equation*}
  p=\text{--suatu bentuk yang melibatkan panjang tali, dan sebagainya--}
\end{equation*}
Periode dinyatakan dalam satuan waktu $L^0M^0T^1$.
Jadi, besaran-besaran pada
ruas lain persamaan harus memiliki rumus dimensi yang berpadu
sedemikian rupa sehingga faktor $L$ dan $M$ saling menghilangkan dan hanya
satu faktor $T$ yang tersisa.
Tabel pada halaman~\pageref{table:Dimen} memuat besaran-besaran yang
oleh seorang peneliti berpengalaman mungkin dianggap relevan terhadap periode
sebuah pendulum.
Satu-satunya rumus dimensi yang melibatkan $L$ adalah rumus untuk
panjang tali dan percepatan gravitasi.
Agar faktor $L$ dari keduanya saling menghilangkan ketika muncul
dalam persamaan, keduanya harus berbentuk rasio,
misalnya sebagai $(\ell/g)^2$, sebagai $\cos(\ell/g)$,
atau sebagai $(\ell/g)^{-1}$.
Oleh karena itu, periode merupakan fungsi dari $\ell/g$.

Ini merupakan hasil yang luar biasa:~hanya dengan analisis menggunakan
pensil dan kertas, sebelum kita mengeluarkan pendulum dan melakukan pengukuran,
kita telah menentukan sesuatu tentang faktor-faktor penyusun periodenya.

Untuk melakukan analisis dimensi secara sistematis, kita memerlukan dua fakta
(argumen bagi kedua fakta ini terdapat dalam \cite{Bridgman}, Bab~II dan~IV).
Fakta pertama adalah bahwa setiap persamaan yang menghubungkan besaran-besaran
fisik yang akan kita jumpai melibatkan jumlah sejumlah suku,
dengan setiap suku berbentuk
\begin{equation*}
  m_1^{p_1}m_2^{p_2}\cdots m_k^{p_k}
\end{equation*}
untuk bilangan $m_1$, \ldots, $m_k$ yang mengukur besaran-besaran tersebut.

Untuk fakta kedua, perhatikan bahwa cara mudah untuk menyusun ekspresi yang
homogen secara dimensi adalah dengan mengambil hasil kali besaran-besaran tak
berdimensi atau menjumlahkan suku-suku tak berdimensi tersebut.
Teorema Buckingham menyatakan bahwa setiap
hubungan lengkap di antara besaran-besaran yang memiliki
rumus dimensi dapat dimanipulasi secara aljabar menjadi bentuk yang memuat
suatu fungsi $f$ sedemikian sehingga
\begin{equation*}
  f(\Pi_1,\ldots,\Pi_n)=0
\end{equation*}
untuk suatu himpunan lengkap $\set{\Pi_1,\ldots,\Pi_n}$ yang terdiri atas
hasil kali tak berdimensi.
(Contoh pertama di bawah menjelaskan apa yang membuat suatu himpunan hasil kali
tak berdimensi disebut `lengkap'.)
Biasanya kita ingin menyatakan salah satu besaran,
misalnya $m_1$, dalam besaran-besaran lainnya.
Untuk itu, kita akan mengasumsikan bahwa persamaan di atas dapat ditulis ulang
\begin{equation*}
  m_1=m_2^{-p_2}\cdots m_k^{-p_k}\cdot \hat{f}(\Pi_2,\ldots,\Pi_n)
\end{equation*}
dengan $\Pi_1=m_1m_2^{p_2}\cdots m_k^{p_k}$ tak berdimensi
dan hasil kali $\Pi_2$, \ldots, $\Pi_n$ tidak melibatkan $m_1$
(seperti $f$, di sini $\hat{f}$ adalah fungsi sembarang, kali ini
dengan $n-1$ argumen).
Jadi, untuk melakukan analisis dimensi, kita harus
menemukan hasil kali tak berdimensi apa saja yang mungkin.
%For that we will use vector spaces.

Sebagai contoh, perhatikan kembali
rumus periode sebuah pendulum.
\begin{center}
  \begin{tabular}{c}
    \includegraphics{vs/mp/ch2.23}
  \end{tabular}
 \qquad\quad
 \label{table:Dimen}
  \begin{tabular}{r|l} 
    \multicolumn{1}{r}{\textit{besaran}}
    &\multicolumn{1}{l}{\begin{tabular}[b]{@{}l} 
                           \textit{rumus} \\[-.3ex]
                           \textit{dimensi}
                         \end{tabular}}                   \\ \hline
    periode $p$                      &$L^0M^0T^1$          \\
    panjang tali $\ell$              &$L^1M^0T^0$          \\
    massa bandul $m$                 &$L^0M^1T^0$          \\
    percepatan gravitasi $g$         &$L^1M^0T^{-2}$    \\
    sudut ayunan $\theta$            &$L^0M^0T^0$
  \end{tabular}
\end{center} 
Menurut fakta pertama yang dikutip di atas, kita mengharapkan rumus tersebut
memiliki bentuk berikut (atau mungkin berupa jumlah suku-suku berbentuk berikut):
$p^{p_1}\ell^{p_2}m^{p_3}g^{p_4}\theta^{p_5}$.
Untuk menggunakan fakta kedua, yakni menemukan kombinasi pangkat
$p_1$, \ldots, $p_5$ yang
menghasilkan hasil kali tak berdimensi, perhatikan persamaan berikut.
\begin{equation*}
  (L^0M^0T^1)^{p_1}(L^1M^0T^0)^{p_2}(L^0M^1T^0)^{p_3}
     (L^1M^0T^{-2})^{p_4}(L^0M^0T^0)^{p_5}
  =L^0M^0T^0 
\end{equation*}
Persamaan itu memberikan tiga syarat bagi pangkat-pangkat tersebut.
\begin{equation*}
  \begin{linsys}{5}
       &\   &p_2  &\   &    &+  &p_4   &  &  &=  &0  \\
       &    &     &    &p_3 &   &      &  &  &=  &0  \\
    p_1&    &     &    &    &-  &2p_4  &  &  &=  &0  
  \end{linsys}  
\end{equation*}
Perhatikan bahwa $p_3=0$, sehingga massa bandul tidak memengaruhi periode.
Reduksi Gauss dan parametrisasi sistem tersebut memberikan
\begin{equation*}
  \set{\colvec{p_1 \\ p_2 \\ p_3 \\ p_4 \\ p_5}=
       \colvec[r]{1 \\  -1/2  \\  0  \\  1/2  \\  0}p_1+
       \colvec[r]{0   \\  0   \\  0  \\  0  \\  1}p_5
       \suchthat p_1,p_5\in\Re}
\end{equation*}
(kita mengambil $p_1$ sebagai salah satu parameter agar dapat menyatakan
periode dalam besaran-besaran lainnya).

Himpunan hasil kali tak berdimensi
memuat semua suku $p^{p_1}\ell^{p_2}m^{p_3}g^{p_4}\theta^{p_5}$ yang memenuhi
syarat-syarat di atas.
Himpunan ini membentuk ruang vektor dengan operasi
`$+$' berupa perkalian dua hasil kali semacam itu
dan operasi `$\cdot$' berupa pemangkatan
hasil kali tersebut dengan skalar
(lihat \nearbyexercise{exer:DimLessProdsVecSp}).
Istilah `himpunan lengkap hasil kali tak berdimensi'
dalam Teorema Buckingham berarti basis bagi ruang vektor ini.

Kita dapat memperoleh suatu basis
dengan mula-mula mengambil $p_1=1$, $p_5=0$, lalu mengambil $p_1=0$, $p_5=1$.
Hasil kali tak berdimensi yang bersesuaian
adalah $\Pi_1=p\ell^{-1/2}g^{1/2}$ dan $\Pi_2=\theta$.
Karena himpunan $\set{\Pi_1,\Pi_2}$ lengkap,
Teorema Buckingham menyatakan bahwa
\begin{equation*}
  p = \ell^{1/2}g^{-1/2}\cdot\hat{f}(\theta) 
    = \sqrt{\ell/g}\cdot\hat{f}(\theta)
\end{equation*}
dengan $\hat{f}$ suatu fungsi yang tidak dapat kita tentukan dari analisis ini
(buku teks fisika tahun pertama akan menunjukkan dengan cara lain
bahwa untuk sudut kecil, fungsi tersebut kira-kira merupakan fungsi konstan
$\hat{f}(\theta)=2\pi$).

Dengan demikian, analisis terhadap hubungan-hubungan yang mungkin di antara
besaran-besaran dengan rumus dimensi yang diberikan
telah memberi kita cukup banyak informasi:~periode pendulum tidak
bergantung pada massa bandul,
dan bertambah sebanding dengan akar kuadrat panjang tali.

Untuk contoh berikutnya, kita mencoba menentukan periode revolusi
dua benda di ruang angkasa yang saling mengorbit akibat tarik-menarik gravitasi.
Seorang peneliti berpengalaman dapat memperkirakan bahwa besaran-besaran
berikutlah yang relevan.
\begin{center}
  \begin{tabular}{@{}c@{}}
    \includegraphics{vs/mp/ch2.24}
  \end{tabular}
  \qquad\quad
  \begin{tabular}{r|l} 
    \multicolumn{1}{r}{\textit{besaran}}
    &\multicolumn{1}{l}{\begin{tabular}[b]{@{}l@{}} 
                          \textit{rumus} \\[-.3ex]
                          \textit{dimensi}
                         \end{tabular}} \\ \hline
    periode $p$                  &$L^0M^0T^1$         \\
    jarak rata-rata $r$          &$L^1M^0T^0$          \\
    massa pertama $m_1$          &$L^0M^1T^0$          \\
    massa kedua $m_2$            &$L^0M^1T^0$          \\
    konstanta gravitasi $G$       &$L^3M^{-1}T^{-2}$
  \end{tabular}
\end{center} 
%(\cite{Bridgman} discusses why the investigator needs experience).
Untuk memperoleh himpunan lengkap hasil kali tak berdimensi, kita meninjau
persamaan
\begin{equation*}
  (L^0M^0T^1)^{p_1}(L^1M^0T^0)^{p_2}(L^0M^1T^0)^{p_3}(L^0M^1T^0)^{p_4}
        (L^3M^{-1}T^{-2})^{p_5}=L^0M^0T^0
\end{equation*}
yang menghasilkan sistem
\begin{equation*}
  \begin{linsys}{5}
          &  &p_2  &   &     &   &      &+  &3p_5 &=  &0  \\
          &  &     &   &p_3  &+  &p_4   &-  &p_5  &=  &0  \\
      p_1 &  &     &   &     &   &      &-  &2p_5 &=  &0  
  \end{linsys}
\end{equation*}
dengan solusi berikut.
\begin{equation*}
  \set{\colvec[r]{1 \\ -3/2  \\ 1/2 \\ 0 \\ 1/2}p_1
       +\colvec[r]{0 \\ 0    \\  -1 \\ 1 \\ 0}p_4
    \suchthat p_1,p_4\in\Re}
\end{equation*}

Seperti sebelumnya,
himpunan hasil kali tak berdimensi dari besaran-besaran ini membentuk ruang vektor,
dan kita ingin menghasilkan suatu basis bagi ruang tersebut,
yakni himpunan `lengkap' hasil kali tak berdimensi.
Salah satu himpunan demikian, yang diperoleh dengan menetapkan $p_1=1$ dan $p_4=0$
serta menetapkan $p_1=0$ dan $p_4=1$, adalah
$\set{\Pi_1=pr^{-3/2}m_1^{1/2}G^{1/2},\,\Pi_2=m_1^{-1}m_2}$.
Dengan itu, Teorema Buckingham menyatakan bahwa
setiap hubungan lengkap di antara besaran-besaran ini dapat dinyatakan dalam bentuk berikut.
\begin{equation*}
  p = r^{3/2}m_1^{-1/2}G^{-1/2}\cdot\hat{f}(m_1^{-1}m_2)  
    = \frac{r^{3/2}}{\sqrt{Gm_1}}\cdot\hat{f}(m_2/m_1)
\end{equation*}

\textit{Catatan.}
Salah satu penerapan penting dari rumus sebelumnya adalah
ketika $m_1$ merupakan massa Matahari dan $m_2$ merupakan massa
sebuah planet.
Karena $m_1$ jauh lebih besar daripada $m_2$,
argumen bagi $\hat{f}$ kira-kira sama dengan $0$,
dan kita dapat bertanya apakah
bagian rumus ini tetap kira-kira konstan ketika $m_2$ berubah.
Salah satu cara untuk melihat bahwa memang demikian adalah sebagai berikut.
Matahari jauh lebih besar daripada planet,
sehingga rotasi timbal balik itu berlangsung kira-kira mengelilingi pusat Matahari.
Jika kita mengubah massa planet $m_2$ dengan faktor $x$
(misalnya, massa Venus adalah
$x=0.815$ kali massa Bumi),
maka gaya tarik dikalikan dengan $x$,
dan gaya sebesar $x$ kali yang bekerja pada massa sebesar $x$ kali,
karena $F=ma$, menghasilkan percepatan yang sama terhadap pusat yang sama
(secara hampiran).
Karena itu,
orbitnya akan sama sehingga periodenya juga sama,
dan dengan demikian ruas kanan persamaan di atas juga tetap tidak berubah
(secara hampiran).
Oleh karena itu,
$\hat{f}(m_2/m_1)$ kira-kira konstan ketika $m_2$ berubah.
Inilah Hukum Ketiga Kepler:~kuadrat
periode sebuah planet sebanding dengan pangkat tiga jari-jari rata-rata orbitnya
mengelilingi Matahari.

Contoh terakhir merupakan salah satu penerapan eksplisit pertama
analisis dimensi.
Lord Rayleigh meninjau kecepatan gelombang di air dalam dan
mengusulkan besaran-besaran berikut sebagai besaran yang relevan.
\begin{center}
  \begin{tabular}{r|l} 
    \multicolumn{1}{r}{\textit{besaran}}
    &\multicolumn{1}{l}{\begin{tabular}[b]{@{}l} 
                          \textit{rumus} \\[-.3ex]
                          \textit{dimensi}
                         \end{tabular}} \\ \hline
    kecepatan gelombang $v$             &$L^1M^0T^{-1}$         \\
    massa jenis air $d$                 &$L^{-3}M^1T^0$         \\
    percepatan gravitasi $g$            &$L^1M^0T^{-2}$         \\
    panjang gelombang $\lambda$         &$L^1M^0T^0$
  \end{tabular}
\end{center} 
Persamaan
\begin{equation*}
  (L^1M^0T^{-1})^{p_1}(L^{-3}M^1T^0)^{p_2}
      (L^1M^0T^{-2})^{p_3}(L^1M^0T^0)^{p_4}=L^0M^0T^0
\end{equation*}
menghasilkan sistem berikut,
\begin{equation*}
  \begin{linsys}{4}
    p_1  &-  &3p_2 &+  &p_3  &+  &p_4  &=  &0  \\
         &   &p_2  &   &     &   &     &=  &0  \\
    -p_1 &   &     &-  &2p_3 &   &     &=  &0  
  \end{linsys}
\end{equation*}
dengan ruang solusi berikut.
\begin{equation*}
  \set{\colvec[r]{1 \\ 0 \\ -1/2  \\ -1/2}p_1
       \suchthat p_1\in \Re}
\end{equation*}
Terdapat satu hasil kali tak berdimensi, $\Pi_1=vg^{-1/2}\lambda^{-1/2}$,
sehingga $v$ sama dengan $\sqrt{\lambda g}$ dikalikan suatu konstanta;
$\hat{f}$ konstan karena merupakan fungsi tanpa argumen.
Besaran $d$
tidak terlibat dalam hubungan tersebut.

Ketiga contoh di atas menunjukkan bahwa analisis dimensi
dapat membawa kita jauh dalam menyatakan hubungan di antara besaran-besaran.
Untuk bacaan lebih lanjut,
rujukan klasiknya adalah \cite{Bridgman}\Dash buku singkat ini sangat menyenangkan.
Sumber lainnya adalah \cite{Giordano83}.
Uraian mengenai kedudukan analisis dimensi dalam
pemodelan terdapat dalam \cite{Giordano82}.

\begin{exercises}
  \item 
    \cite{deMestre}
    Perhatikan sebuah proyektil yang diluncurkan dengan kecepatan awal $v_0$
    pada sudut $\theta$.
    Untuk mempelajari geraknya, kita dapat
    menduga bahwa besaran-besaran berikutlah yang relevan.
    \begin{center}
      \begin{tabular}{r|l} 
        \multicolumn{1}{r}{\textit{besaran}}
        &\multicolumn{1}{l}{\begin{tabular}[b]{@{}l} 
                               \textit{rumus} \\[-.3ex]
                               \textit{dimensi}
                             \end{tabular}} \\ \hline
        posisi horizontal $x$            &$L^1M^0T^0$         \\
        posisi vertikal $y$              &$L^1M^0T^0$         \\
        kelajuan awal $v_0$              &$L^1M^0T^{-1}$      \\
        sudut peluncuran $\theta$         &$L^0M^0T^0$         \\
        percepatan gravitasi $g$          &$L^1M^0T^{-2}$      \\
        waktu $t$                         &$L^0M^0T^1$
      \end{tabular}
    \end{center} 
    \begin{exparts}
      \partsitem Tunjukkan bahwa $\set{gt/v_0,gx/v_0^2,gy/v_0^2,\theta}$
        merupakan himpunan lengkap hasil kali tak berdimensi.
        (\textit{Petunjuk.}
        Salah satu caranya adalah menemukan variabel bebas yang sesuai
        dalam sistem linear yang muncul, tetapi ada jalan pintas yang
        menggunakan sifat-sifat basis.)
      \partsitem Kedua persamaan gerak proyektil berikut sudah dikenal:
        $x=v_0\cos(\theta) t$ dan $y=v_0\sin(\theta) t- (g/2)t^2$.
        Manipulasilah masing-masing persamaan untuk menuliskannya ulang sebagai
        hubungan di antara hasil kali tak berdimensi pada butir sebelumnya.
    \end{exparts}
    \begin{answer}
      \begin{exparts}
        \partsitem Dengan mengasumsikan bahwa
          \begin{equation*}
            (L^1M^0T^0)^{p_1}(L^1M^0T^0)^{p_2}(L^1M^0T^{-1})^{p_3}
             (L^0M^0T^0)^{p_4}(L^1M^0T^{-2})^{p_5}(L^0M^0T^1)^{p_6}
          \end{equation*}
          sama dengan $L^0M^0T^0$, kita memperoleh sistem linear berikut,
          \begin{equation*}
            \begin{linsys}{5}
              p_1  &+  &p_2  &+  &p_3  &+  &p_5  &   &    &=  &0  \\
              %      &   &     &   &     &   &     &   &0   &=  &0  \\
                   &   &     &   &-p_3 &-  &2p_5 &+  &p_6 &=  &0  
            \end{linsys}
          \end{equation*}
          (tidak ada pembatasan pada $p_4$).
          Parametrisasi yang alami menggunakan variabel bebas untuk memberikan
          $p_3=-2p_5+p_6$ dan $p_1=-p_2+p_5-p_6$.
          Uraian himpunan solusi yang dihasilkan,
          \begin{equation*}
            \set{\colvec{p_1 \\ p_2 \\ p_3 \\ p_4 \\ p_5 \\ p_6}
                =p_2\colvec[r]{-1 \\ 1 \\ 0 \\ 0  \\ 0 \\ 0} 
                +p_4\colvec[r]{0 \\ 0 \\ 0 \\ 1 \\ 0 \\ 0} 
                +p_5\colvec[r]{1 \\ 0 \\ -2 \\ 0 \\ 1 \\ 0} 
                +p_6\colvec[r]{-1 \\ 0 \\ 1 \\ 0 \\ 0 \\ 1} 
                \suchthat p_2, p_4, p_5, p_6\in\Re  }
          \end{equation*}
          memberikan $\set{y/x,\theta,xt/{v_0}^2,v_0t/x}$
          sebagai himpunan lengkap hasil kali tak berdimensi
          (ingat bahwa ``lengkap'' dalam konteks ini tidak berarti
          tidak ada hasil kali tak berdimensi lainnya;
          istilah itu hanya berarti bahwa himpunan tersebut adalah basis).
          Namun, ini bukan himpunan hasil kali tak berdimensi yang
          diminta dalam soal.

          Ada dua cara untuk melanjutkan.
          Cara pertama adalah mengutak-atik pilihan parameter dengan harapan
          memperoleh himpunan yang tepat.
          Untuk itu, kita dapat menjalankan langkah pada paragraf sebelumnya secara terbalik.
          Mengubah hasil kali tak berdimensi yang diberikan, yaitu $gt/v_0$,
          $gx/v_0^2$, $gy/v_0^2$, dan $\theta$
          menjadi vektor memberikan uraian berikut (perhatikan tanda \textit{?}
          pada tempat parameter akan diletakkan).
          \begin{equation*}
            \set{\colvec{p_1 \\ p_2 \\ p_3 \\ p_4 \\ p_5 \\ p_6}
                =\underline{\ \text{\textit{?}}\ }
                      \colvec[r]{0  \\ 0 \\ -1 \\ 0  \\ 1 \\ 1} 
                +\underline{\ \text{\textit{?}}\ }
                   \colvec[r]{1 \\ 0 \\ -2 \\ 0 \\ 1 \\ 0} 
                +\underline{\ \text{\textit{?}}\ }
                   \colvec[r]{0 \\ 1 \\ -2 \\ 0 \\ 1 \\ 0} 
                +p_4\colvec[r]{0  \\ 0 \\ 0 \\ 1 \\ 0 \\ 0} 
                \suchthat p_2, p_4, p_5, p_6\in\Re  }
          \end{equation*}
          $p_4$ sudah berada pada tempatnya.
          Pemeriksaan terhadap baris-baris menunjukkan bahwa kita juga dapat
          menempatkan $p_6$, $p_1$, dan $p_2$.

          Cara kedua, dengan mengikuti petunjuk, adalah memperhatikan bahwa
          himpunan yang diberikan berukuran empat dalam ruang vektor berdimensi empat,
          sehingga kita hanya perlu menunjukkan bahwa himpunan itu bebas linear.
          Hal ini mudah dilakukan melalui pemeriksaan, dengan meninjau
          komponen keenam, pertama, kedua, dan keempat dari vektor-vektor tersebut.
        \item Persamaan pertama dapat ditulis ulang sebagai
          \begin{equation*}
            \frac{gx}{{v_0}^2}=\frac{gt}{v_0}\cos\theta
          \end{equation*}
          sehingga fungsi Buckingham adalah
          $f_1(\Pi_1,\Pi_2,\Pi_3,\Pi_4)=\Pi_2-\Pi_1\cos(\Pi_4)$.
          Persamaan kedua dapat ditulis ulang sebagai
          \begin{equation*}
            \frac{gy}{{v_0}^2}=\frac{gt}{v_0}\sin\theta
              -\frac{1}{2}\left(\frac{gt}{v_0}\right)^2
          \end{equation*}
          dan fungsi Buckingham di sini adalah
          $f_2(\Pi_1,\Pi_2,\Pi_3,\Pi_4)
             =\Pi_3-\Pi_1\sin(\Pi_4)+(1/2){\Pi_1}^2$.
      \end{exparts}
    \end{answer}
  \item \cite{Einstein1911}
    menduga bahwa frekuensi karakteristik inframerah suatu benda padat
    mungkin ditentukan oleh gaya antaratom yang sama dengan gaya yang menentukan
    perilaku elastis biasa benda padat tersebut.
    Besaran-besaran yang relevan adalah sebagai berikut.
    \begin{center}
      \begin{tabular}{r|l} 
        \multicolumn{1}{r}{\textit{besaran}}
        &\multicolumn{1}{l}{\begin{tabular}[b]{@{}l} 
                              \textit{rumus} \\[-.3ex]
                              \textit{dimensi}
                             \end{tabular}} \\ \hline
        frekuensi karakteristik $\nu$         &$L^0M^0T^{-1}$         \\
        kompresibilitas $k$                   &$L^1M^{-1}T^2$          \\
        jumlah atom per cm kubik $N$          &$L^{-3}M^0T^0$          \\
        massa sebuah atom $m$                 &$L^0M^1T^0$
      \end{tabular}
    \end{center} 
    Tunjukkan bahwa terdapat satu hasil kali tak berdimensi.
    Simpulkan bahwa dalam setiap hubungan lengkap di antara besaran-besaran
    dengan rumus dimensi tersebut, $k$ sama dengan suatu konstanta dikalikan
    $\nu^{-2}N^{-1/3}m^{-1}$.
    Kesimpulan ini memainkan peran penting dalam kajian awal
    fenomena kuantum.
    \begin{answer}
      Perhatikan
      \begin{equation*}
        (L^0M^0T^{-1})^{p_1}(L^1M^{-1}T^2)^{p_2}
          (L^{-3}M^0T^0)^{p_3}(L^0M^1T^0)^{p_4}=(L^0M^0T^0)
      \end{equation*}
      yang memberikan hubungan-hubungan berikut di antara pangkat-pangkat tersebut.
      \begin{equation*}
        \begin{linsys}{4}
               &    &p_2   &-  &3p_3  &   &    &=  &0  \\
               &    &-p_2  &   &      &+  &p_4 &=  &0  \\
         -p_1  &+   &2p_2  &   &      &   &    &=  &0  
        \end{linsys}
        \grstep{\rho_1\leftrightarrow\rho_3}
        \repeatedgrstep{\rho_2+\rho_3}
        \begin{linsys}{4}
         -p_1  &+   &2p_2  &   &      &   &    &=  &0  \\
               &    &-p_2  &   &      &+  &p_4 &=  &0  \\
               &    &      &\  &-3p_3 &+  &p_4 &=  &0  
        \end{linsys}
      \end{equation*}
      Berikut adalah ruang solusinya
      (karena kita ingin menyatakan $k$ sebagai fungsi besaran-besaran lainnya,
      kita mengambil $p_2$ sebagai parameter).
      \begin{equation*}
        \set{\colvec[r]{2 \\ 1 \\ 1/3 \\ 1}p_2
             \suchthat p_2\in\Re}
      \end{equation*}
      Dengan demikian, $\Pi_1=\nu^2kN^{1/3}m$ merupakan kombinasi tak berdimensi,
      dan kita memperoleh bahwa $k$ sama dengan
      $\nu^{-2}N^{-1/3}m^{-1}$ dikalikan suatu konstanta
      (fungsi $\hat{f}$ konstan karena tidak memiliki argumen).
    \end{answer}
  \item 
    \cite{Giordano83}
    Torsi yang dihasilkan oleh sebuah mesin
    memiliki rumus dimensi $L^2M^1T^{-2}$.
    Mula-mula kita dapat menduga bahwa torsi tersebut bergantung pada
    laju putaran mesin (dengan rumus dimensi
    $L^0M^0T^{-1}$) dan volume udara yang dipindahkan (dengan rumus
    dimensi $L^3M^0T^0$).
    \begin{exparts}
      \partsitem Cobalah mencari himpunan lengkap hasil kali tak berdimensi.
        Apa yang tidak berjalan semestinya?
      \partsitem Sesuaikan dugaan tersebut dengan menambahkan massa jenis udara
        (dengan rumus dimensi $L^{-3}M^1T^0$).
        Sekarang carilah himpunan lengkap hasil kali tak berdimensi.
    \end{exparts}
    \begin{answer}
      \begin{exparts}
        \partsitem Menetapkan
          \begin{equation*}
            (L^2M^1T^{-2})^{p_1}(L^0M^0T^{-1})^{p_2}(L^{3}M^0T^0)^{p_3}
              =(L^0M^0T^0)
          \end{equation*}
          menghasilkan
          \begin{equation*}
            \begin{linsys}{3}
              2p_1  &   &    &+  &3p_3  &=  &0  \\
              p_1   &   &    &   &      &=  &0  \\
              -2p_1 &-  &p_2 &   &      &=  &0  
            \end{linsys}
          \end{equation*}
          yang menyiratkan bahwa $p_1=p_2=p_3=0$.
          Artinya, di antara besaran-besaran dengan rumus dimensi ini, satu-satunya
          hasil kali tak berdimensi adalah hasil kali trivial.
        \partsitem Menetapkan
          \begin{equation*}
            (L^2M^1T^{-2})^{p_1}(L^0M^0T^{-1})^{p_2}(L^{3}M^0T^0)^{p_3}
              (L^{-3}M^1T^0)^{p_4}=(L^0M^0T^0)
          \end{equation*}
          menghasilkan sistem berikut.
          \begin{multline*}
            \begin{linsys}{4}
              2p_1  &   &    &+  &3p_3  &-  &3p_4 &=  &0  \\
              p_1   &   &    &   &      &+  &p_4  &=  &0  \\
              -2p_1 &-  &p_2 &   &      &   &     &=  &0  
            \end{linsys}                                         \\
            \grstep[\rho_1+\rho_3]{(-1/2)\rho_1+\rho_2}
            \repeatedgrstep{\rho_2\leftrightarrow\rho_3}
            \begin{linsys}{4}
              2p_1  &   &      &+  &3p_3       &-   &3p_4      &=  &0  \\
                    &   &-p_2  &+  &3p_3       &-   &3p_4      &=  &0  \\
                    &  &      &   &(-3/2)p_3  &+   &(5/2)p_4  &=  &0  
            \end{linsys}
          \end{multline*}
          Mengambil $p_1$ sebagai parameter untuk menyatakan torsi memberikan
          uraian himpunan solusi berikut.
          \begin{equation*}
            \set{\colvec[r]{1 \\ -2 \\ -5/3 \\ -1}p_1
                 \suchthat p_1\in\Re}
          \end{equation*}
          Dengan menyatakan torsi sebagai $\tau$, laju putaran sebagai $r$, volume
          udara sebagai $V$, dan massa jenis udara sebagai $d$, kita memperoleh
          $\Pi_1=\tau r^{-2} V^{-5/3} d^{-1}$, sehingga
          torsi sama dengan $r^2V^{5/3}d$ dikalikan suatu konstanta.
      \end{exparts}
    \end{answer}
  \item 
    \cite{Tilley} 
    Domino-domino yang jatuh membentuk gelombang.
    Kita dapat menduga bahwa kecepatan gelombang $v$ bergantung pada
    jarak $d$ di antara domino-domino tersebut,
    tinggi $h$ setiap domino,
    dan percepatan gravitasi $g$.
    \begin{exparts}
      \partsitem Carilah rumus dimensi untuk masing-masing dari keempat besaran.
      \partsitem Tunjukkan bahwa $\set{\Pi_1=h/d,\Pi_2=dg/v^2}$
        merupakan himpunan lengkap hasil kali tak berdimensi.
      \partsitem Tunjukkan bahwa jika $h/d$ dibuat tetap, maka
        kecepatan rambat sebanding dengan akar kuadrat
        dari $d$.
    \end{exparts}
    \begin{answer}
          \begin{exparts}
            \partsitem Berikut adalah rumus-rumus dimensinya.
              \begin{center}
                \begin{tabular}{r|l} 
                  \multicolumn{1}{r}{\textit{besaran}}
                    &\multicolumn{1}{l}{\begin{tabular}[b]{@{}l} 
                                           \textit{rumus} \\[-.3ex]
                                           \textit{dimensi}
                                       \end{tabular}} \\ \hline
                        kecepatan gelombang $v$             &$L^1M^0T^{-1}$ \\
                        jarak antar-domino $d$               &$L^1M^0T^0$   \\
                        tinggi domino $h$                    &$L^1M^0T^0$ \\
                        percepatan gravitasi $g$             &$L^1M^0T^{-2}$
                \end{tabular}
              \end{center}     
            \partsitem Hubungan
              \begin{equation*}
                (L^1M^0T^{-1})^{p_1}(L^1M^0T^0)^{p_2}(L^1M^0T^0)^{p_3}
                  (L^1M^0T^{-2})^{p_4}=(L^0M^0T^0)
              \end{equation*}
              menghasilkan sistem linear berikut.
              \begin{equation*}
                \begin{linsys}{4}
                  p_1  &+  &p_2  &+  &p_3  &+  &p_4  &=  &0  \\
                       &   &     &   &     &   &0    &=  &0  \\
                 -p_1  &   &     &   &     &-  &2p_4 &=  &0  
                \end{linsys}
                \grstep{\rho_1+\rho_4}
                \begin{linsys}{4}
                  p_1  &+  &p_2  &+  &p_3  &+  &p_4  &=  &0  \\
                       &   &p_2  &+  &p_3  &-  &p_4  &=  &0  
                \end{linsys}
              \end{equation*}
              Dengan mengambil $p_3$ dan $p_4$ sebagai parameter, kita dapat
              menguraikan himpunan solusi sebagai berikut.
              \begin{equation*}
                \set{\colvec[r]{0 \\ -1 \\ 1 \\ 0}p_3
                     +\colvec[r]{-2 \\ 1 \\ 0 \\ 1}p_4
                     \suchthat p_3,p_4\in\Re}
              \end{equation*}
              Hasilnya memberikan $\set{\Pi_1=h/d,\Pi_2=dg/v^2}$ sebagai
              suatu himpunan lengkap.
            \partsitem Teorema Buckingham menyatakan bahwa
              $v^2=dg\cdot\hat{f}(h/d)$, sehingga karena $g$ adalah konstanta,
              jika $h/d$ dibuat tetap, maka $v$ sebanding dengan $\sqrt{d\,}$.
          \end{exparts}
        \end{answer}
   \item \label{exer:DimLessProdsVecSp}
      Buktikan bahwa hasil kali tak berdimensi membentuk ruang vektor
      dengan operasi $\vec{+}$ berupa perkalian dua hasil kali semacam itu
      dan operasi $\vec{\cdot}$ berupa pemangkatan
      hasil kali tersebut dengan skalar.
      (Panah vektor digunakan untuk mencegah kerancuan.)
      Artinya, buktikan bahwa untuk setiap sistem homogen tertentu,
      himpunan hasil kali pangkat-pangkat $m_1$, \ldots, $m_k$ berikut,
      \begin{equation*}
        \set{m_1^{p_1}\dots m_k^{p_k}
             \suchthat 
              \text{$p_1$, \ldots, $p_k$ memenuhi sistem tersebut}}
      \end{equation*}
      merupakan ruang vektor dengan operasi berikut:
      \begin{equation*}
        m_1^{p_1}\dots m_k^{p_k}
        \vec{+}
        m_1^{q_1}\dots m_k^{q_k}
        =
        m_1^{p_1+q_1}\dots m_k^{p_k+q_k}
      \end{equation*}
      dan
      \begin{equation*}
        r\vec{\cdot} (m_1^{p_1}\dots m_k^{p_k})
        =
        m_1^{rp_1}\dots m_k^{rp_k}
      \end{equation*}
      (asumsikan bahwa semua variabel menyatakan bilangan real).
      \begin{answer}
        Pemeriksaan syarat-syarat dalam definisi ruang vektor bersifat
        rutin.
      \end{answer}
%  \item Buckingham's Theorem says there exists a complete set such that
%    the relationship can be stated in that way.
%    Why do we not have to worry that the particular complete set
%    we have produced might not be the right one?
%    (\textit{Hint.} 
%    Consider the prior exercise.)
%    \begin{answer}
%      Any `complete' set (any basis for the vector space under consideration)
%      can be expressed in terms of any other one.
%      So any function that is correct when given the members of a particular
%      complete set as arguments can
%      be rewritten to be correct for the  members of any other complete set.
%    \end{answer}
  \item \label{exer:AppsAndOranges}
    Nasihat tentang apel dan jeruk itu tidak sepenuhnya benar.
    Perhatikan persamaan lingkaran yang sudah dikenal, $C=2\pi r$ dan $A=\pi r^2$.
    \begin{exparts}
      \partsitem Periksa bahwa $C$ dan $A$ memiliki rumus dimensi yang berbeda.
      \partsitem Susunlah sebuah persamaan yang tidak
        homogen secara dimensi (yakni,
        menjumlahkan apel dan jeruk), tetapi tetap benar untuk setiap lingkaran.
      \partsitem Butir sebelumnya meminta persamaan yang lengkap tetapi
        tidak homogen secara dimensi.
        Susunlah persamaan yang homogen secara dimensi tetapi tidak lengkap.
    \end{exparts}
    (Walaupun pepatah lama itu tidak sepenuhnya benar, pepatah tersebut tetap
    memberikan strategi yang berguna.
    Kehomogenan dimensi sering digunakan untuk memeriksa kewajaran
    persamaan yang digunakan dalam model.
    Untuk argumen bahwa setiap persamaan lengkap dapat dengan mudah dibuat
    homogen secara dimensi, lihat \cite{Bridgman}, Bab~I,
    khususnya halaman~15.)
    \begin{answer}
      \begin{exparts}
        \partsitem Rumus dimensi keliling adalah $L$,
          yakni $L^1M^0T^0$.
          Rumus dimensi luas adalah $L^2$.
        \partsitem Salah satunya adalah $C+A=2\pi r + \pi r^2$.
        \partsitem Salah satu contohnya adalah rumus berikut yang menghubungkan
           panjang busur yang dibentuk oleh suatu sudut dengan jari-jari dan
           ukuran sudut dalam radian:~$\ell-r\theta=0$.
           Kedua suku dalam rumus tersebut memiliki
           rumus dimensi $L^1$.
           Hubungan itu berlaku untuk beberapa
           sistem satuan (misalnya inci dan radian), tetapi tidak untuk semua
           sistem satuan (misalnya inci dan derajat).
      \end{exparts}
    \end{answer}
\end{exercises}
\endinput
